\subsubsection{Vzájemná poloha přímek zadaných v různých tvarech}
\label{subsubsec:Relative_position_mixed_line_forms}

V předchozích částech byly obě přímky zadány stejným způsobem.
V úlohách však mohou být jednotlivé přímky zadány také různými tvary.

Jedna přímka může být například zadána parametricky a druhá obecnou
rovnicí. V takovém případě můžeme jednu z nich převést do stejného
tvaru jako druhou pomocí postupů uvedených v podsekci
\ref{subsec:Conversions_between_line_forms}. Úplný převod však není
vždy nutný. Často postačí z obou zápisů určit údaj popisující směr
přímky a tyto údaje porovnat.

Při určování vzájemné polohy budeme používat především následující
informace:

\begingroup
\renewcommand{\arraystretch}{1.6}
\[
\begin{array}{c|c|c}
    \text{způsob zadání}
    &
    \text{údaj určující směr}
    &
    \text{další užitečný údaj}
    \\
    \hline
    \text{parametrické vyjádření}
    &
    \vec{u}=(u_x;u_y)
    &
    A[x_A;y_A]
    \\[1mm]
    \text{obecná rovnice}
    &
    \vec{n}=(a;b)
    &
    \vec{u}=(-b;a)
    \\[1mm]
    \text{směrnicový tvar}
    &
    k
    &
    \vec{u}=(1;k),\quad Q[0;q]
\end{array}
\]
\endgroup

Při řešení tedy můžeme postupovat dvěma způsoby:

\begin{enumerate}
    \item jednu přímku převedeme do stejného tvaru jako druhou
    a použijeme již známý postup;
    \item bez úplného převodu porovnáme směrové nebo normálové vektory,
    případně směrnice, a průsečík určíme přímým dosazením jednoho
    zápisu do druhého.
\end{enumerate}

Druhý postup bývá zpravidla kratší, protože využívá pouze ty údaje,
které jsou pro určení vzájemné polohy skutečně potřebné.


\paragraph{Parametrické vyjádření a obecná rovnice}

Uvažujme přímku $p$ zadanou parametricky
\begin{equation}\label{eq:Mixed_parametric_general_p}
    \begin{aligned}
        p:\quad x&=x_A+t u_x,\\
                 y&=y_A+t u_y,
                 \qquad t\in\mathbb{R},
    \end{aligned}
\end{equation}
a přímku $q$ zadanou obecnou rovnicí
\begin{equation}\label{eq:Mixed_parametric_general_q}
    q:\quad ax+by+c=0.
\end{equation}

Směr přímky $p$ je určen směrovým vektorem
\[
    \vec{u}_p=(u_x;u_y),
\]
zatímco přímka $q$ má normálový vektor
\[
    \vec{n}_q=(a;b).
\]

Přímky mají stejný směr právě tehdy, když je směrový vektor přímky
$p$ kolmý k normálovému vektoru přímky $q$. Tuto podmínku můžeme
zapsat jako
\begin{equation}\label{eq:Mixed_parametric_general_parallel_condition}
    \vec{u}_p\cdot\vec{n}_q
    =
    au_x+bu_y
    =
    0.
\end{equation}

Jestliže
\[
    au_x+bu_y\neq0,
\]
přímky nemají stejný směr, a jsou proto různoběžné.

Jestliže naopak
\[
    au_x+bu_y=0,
\]
mají přímky stejný směr. Potom musíme ještě rozhodnout, zda jsou
rovnoběžné, nebo totožné. Stačí ověřit, zda zvolený bod
$A[x_A;y_A]$ parametricky zadané přímky leží také na přímce $q$:

\[
\begin{array}{c|c}
    ax_A+by_A+c\neq0
    &
    p\parallel q
    \\[1mm]
    ax_A+by_A+c=0
    &
    p=q
\end{array}
\]

Průsečík různoběžných přímek můžeme určit bez převodu celé rovnice.
Parametrické rovnice přímky $p$ dosadíme přímo do obecné rovnice
přímky $q$:
\begin{align*}
    a(x_A+t u_x)+b(y_A+t u_y)+c&=0,\\
    ax_A+at u_x+by_A+bt u_y+c&=0,\\
    t(au_x+bu_y)+ax_A+by_A+c&=0.
\end{align*}

Je-li $au_x+bu_y\neq0$, získáme jedinou hodnotu parametru
\begin{equation}\label{eq:Mixed_parametric_general_intersection_parameter}
    t
    =
    -\frac{ax_A+by_A+c}{au_x+bu_y}.
\end{equation}
Po jejím dosazení do parametrického vyjádření přímky $p$ získáme
souřadnice průsečíku.

\begin{remark}\label{rem:Mixed_parametric_general_equation_meaning}
    Rovnice vzniklá dosazením parametrického vyjádření do obecné
    rovnice může mít tři různé podoby:

    \begin{itemize}
        \item má právě jedno řešení pro $t$; přímky jsou různoběžné;
        \item vznikne nepravdivý výrok, například $5=0$; přímky jsou
        rovnoběžné a různé;
        \item vznikne pravdivý výrok, například $0=0$; každý bod
        parametricky zadané přímky splňuje obecnou rovnici, a přímky
        jsou proto totožné.
    \end{itemize}
\end{remark}

Kolmost lze určit také přímo. Přímka $q$ je kolmá ke svému normálovému
vektoru $\vec{n}_q$. Proto platí
\[
    p\perp q
\]
právě tehdy, když jsou vektory $\vec{u}_p$ a $\vec{n}_q$ rovnoběžné.

\begin{example}\label{exa:Mixed_parametric_general_intersection}
    Určete vzájemnou polohu přímek
    \[
        \begin{aligned}
            p:\quad x&=1+2t,\\
                     y&=2-t,
                     \qquad t\in\mathbb{R},
        \end{aligned}
    \]
    a
    \[
        q:\quad x+y-4=0.
    \]
    Jsou-li přímky různoběžné, určete jejich průsečík.

    Směrový vektor přímky $p$ a normálový vektor přímky $q$ jsou
    \[
        \vec{u}_p=(2;-1),
        \qquad
        \vec{n}_q=(1;1).
    \]

    Vypočítáme jejich skalární součin:
    \[
        \vec{u}_p\cdot\vec{n}_q
        =
        (2;-1)\cdot(1;1)
        =
        2-1
        =
        1.
    \]

    Skalární součin není roven nule, takže směrový vektor přímky $p$
    není kolmý k normálovému vektoru přímky $q$. Přímky tedy nemají
    stejný směr a jsou různoběžné.

    Průsečík určíme dosazením parametrických rovnic přímky $p$ do
    obecné rovnice přímky $q$:
    \begin{align*}
        (1+2t)+(2-t)-4&=0,\\
        1+2t+2-t-4&=0,\\
        t-1&=0,\\
        t&=1.
    \end{align*}

    Hodnotu $t=1$ dosadíme do parametrického vyjádření:
    \begin{align*}
        x&=1+2\cdot1=3,\\
        y&=2-1=1.
    \end{align*}

    Průsečíkem přímek je bod
    \[
        P[3;1].
    \]

    Výsledek ověříme dosazením do obecné rovnice:
    \[
        3+1-4=0.
    \]

    Platí tedy
    \[
        p\cap q=\{P[3;1]\}.
    \]
\end{example}


\paragraph{Parametrické vyjádření a směrnicový tvar}

Uvažujme přímku $p$ zadanou parametricky
\begin{equation}\label{eq:Mixed_parametric_slope_p}
    \begin{aligned}
        p:\quad x&=x_A+t u_x,\\
                 y&=y_A+t u_y,
                 \qquad t\in\mathbb{R},
    \end{aligned}
\end{equation}
a přímku $q$ ve směrnicovém tvaru
\begin{equation}\label{eq:Mixed_parametric_slope_q}
    q:\quad y=kx+q_0.
\end{equation}

Směrový vektor přímky $p$ je
\[
    \vec{u}_p=(u_x;u_y),
\]
zatímco jako směrový vektor přímky $q$ můžeme zvolit
\[
    \vec{u}_q=(1;k).
\]

Přímky mají stejný směr právě tehdy, když jsou jejich směrové vektory
rovnoběžné. Jestliže $u_x\neq0$, můžeme tuto podmínku zapsat jako
\[
    \frac{u_y}{u_x}=k.
\]

Podmínku lze zapsat také bez dělení:
\begin{equation}\label{eq:Mixed_parametric_slope_parallel_condition}
    u_y-ku_x=0.
\end{equation}

Tento zápis je výhodný také při rozlišování svislé přímky. Je-li
$u_x=0$, potom musí platit $u_y\neq0$, a proto
\[
    u_y-ku_x=u_y\neq0.
\]
Svislá parametricky zadaná přímka tedy nemůže být rovnoběžná s přímkou
ve směrnicovém tvaru a vždy ji protíná.

Jestliže
\[
    u_y-ku_x\neq0,
\]
jsou přímky různoběžné.

Jestliže
\[
    u_y-ku_x=0,
\]
mají přímky stejný směr. Jejich totožnost ověříme dosazením bodu
$A[x_A;y_A]$ do směrnicové rovnice:
\[
\begin{array}{c|c}
    y_A\neq kx_A+q_0
    &
    p\parallel q
    \\[1mm]
    y_A=kx_A+q_0
    &
    p=q
\end{array}
\]

Průsečík různoběžných přímek určíme dosazením parametrických rovnic
do směrnicového tvaru:
\begin{align*}
    y_A+t u_y
    &=
    k(x_A+t u_x)+q_0,\\
    y_A+t u_y
    &=
    kx_A+kt u_x+q_0,\\
    t(u_y-ku_x)
    &=
    kx_A+q_0-y_A.
\end{align*}

Je-li $u_y-ku_x\neq0$, dostáváme
\begin{equation}\label{eq:Mixed_parametric_slope_intersection_parameter}
    t
    =
    \frac{kx_A+q_0-y_A}{u_y-ku_x}.
\end{equation}

Kolmost přímek můžeme určit pomocí skalárního součinu jejich směrových
vektorů:
\[
    \vec{u}_p\cdot\vec{u}_q
    =
    (u_x;u_y)\cdot(1;k)
    =
    u_x+ku_y.
\]
Proto platí
\begin{equation}\label{eq:Mixed_parametric_slope_perpendicular_condition}
    p\perp q
    \quad\Longleftrightarrow\quad
    u_x+ku_y=0.
\end{equation}

\begin{example}\label{exa:Mixed_parametric_slope_parallel_coincident}
    Je dána parametricky zadaná přímka
    \[
        \begin{aligned}
            p:\quad x&=1+2t,\\
                     y&=2+t,
                     \qquad t\in\mathbb{R},
        \end{aligned}
    \]
    a přímky ve směrnicovém tvaru
    \[
        q:\quad y=\frac12x+3,
    \]
    \[
        r:\quad y=\frac12x+\frac32.
    \]

    Určete vzájemnou polohu dvojic přímek $p$, $q$ a přímek $p$, $r$.

    Směrový vektor přímky $p$ je
    \[
        \vec{u}_p=(2;1),
    \]
    a její směrnice je proto
    \[
        k_p=\frac12.
    \]

    Přímky $q$ i $r$ mají rovněž směrnici
    \[
        k_q=k_r=\frac12.
    \]
    Všechny tři přímky tedy mají stejný směr.

    Nejprve ověříme, zda bod $A[1;2]$ přímky $p$ leží na přímce $q$:
    \[
        \frac12\cdot1+3
        =
        \frac72
        \neq
        2.
    \]
    Bod $A$ na přímce $q$ neleží. Přímky $p$ a $q$ mají stejný směr,
    ale nejsou totožné, a proto
    \[
        p\parallel q.
    \]

    Nyní dosadíme bod $A[1;2]$ do rovnice přímky $r$:
    \[
        \frac12\cdot1+\frac32
        =
        2.
    \]
    Bod $A$ na přímce $r$ leží. Přímky $p$ a $r$ mají stejný směr
    a společný bod, a proto jsou totožné:
    \[
        p=r.
    \]
\end{example}


\paragraph{Obecná rovnice a směrnicový tvar}

Uvažujme přímku $p$ zadanou obecnou rovnicí
\begin{equation}\label{eq:Mixed_general_slope_p}
    p:\quad ax+by+c=0
\end{equation}
a přímku $q$ ve směrnicovém tvaru
\begin{equation}\label{eq:Mixed_general_slope_q}
    q:\quad y=kx+q_0.
\end{equation}

Normálový vektor přímky $p$ je
\[
    \vec{n}_p=(a;b),
\]
zatímco směrový vektor přímky $q$ můžeme zvolit jako
\[
    \vec{u}_q=(1;k).
\]

Přímky mají stejný směr právě tehdy, když je směrový vektor přímky
$q$ kolmý k normálovému vektoru přímky $p$:
\[
    \vec{n}_p\cdot\vec{u}_q
    =
    (a;b)\cdot(1;k)
    =
    a+bk.
\]

Podmínka stejného směru tedy je
\begin{equation}\label{eq:Mixed_general_slope_parallel_condition}
    a+bk=0.
\end{equation}

Je-li $b=0$, obecná rovnice popisuje svislou přímku. Protože potom
$a\neq0$, platí
\[
    a+bk=a\neq0.
\]
Svislá přímka se tedy s každou přímkou ve směrnicovém tvaru protíná.

Jestliže
\[
    a+bk\neq0,
\]
jsou přímky různoběžné.

Jestliže
\[
    a+bk=0,
\]
mají přímky stejný směr. K rozlišení rovnoběžnosti a totožnosti
použijeme například bod
\[
    Q[0;q_0],
\]
který leží na přímce $q$. Po dosazení do obecné rovnice přímky $p$
dostáváme
\[
\begin{array}{c|c}
    bq_0+c\neq0
    &
    p\parallel q
    \\[1mm]
    bq_0+c=0
    &
    p=q
\end{array}
\]

Průsečík různoběžných přímek získáme dosazením směrnicové rovnice do
obecné rovnice:
\begin{align*}
    ax+b(kx+q_0)+c&=0,\\
    ax+bkx+bq_0+c&=0,\\
    (a+bk)x+bq_0+c&=0.
\end{align*}

Je-li $a+bk\neq0$, dostáváme
\begin{equation}\label{eq:Mixed_general_slope_intersection_x}
    x
    =
    -\frac{bq_0+c}{a+bk}.
\end{equation}
Druhou souřadnici průsečíku určíme ze vztahu
\[
    y=kx+q_0.
\]

Kolmost můžeme určit porovnáním směrových vektorů. Jako směrový vektor
přímky $p$ můžeme zvolit
\[
    \vec{u}_p=(-b;a),
\]
zatímco pro přímku $q$ máme
\[
    \vec{u}_q=(1;k).
\]
Přímky jsou kolmé právě tehdy, když
\[
    \vec{u}_p\cdot\vec{u}_q
    =
    (-b;a)\cdot(1;k)
    =
    -b+ak
    =
    0.
\]
Tedy
\begin{equation}\label{eq:Mixed_general_slope_perpendicular_condition}
    p\perp q
    \quad\Longleftrightarrow\quad
    ak-b=0.
\end{equation}

\begin{example}\label{exa:Mixed_general_slope_perpendicular}
    Určete vzájemnou polohu přímek
    \[
        p:\quad x+2y-5=0
    \]
    a
    \[
        q:\quad y=2x.
    \]
    Jsou-li přímky různoběžné, určete jejich průsečík.

    Pro obecnou rovnici přímky $p$ platí
    \[
        a=1,
        \qquad
        b=2,
        \qquad
        c=-5.
    \]
    Směrnice přímky $q$ je
    \[
        k=2,
        \qquad
        q_0=0.
    \]

    Nejprve ověříme podmínku stejného směru:
    \[
        a+bk
        =
        1+2\cdot2
        =
        5.
    \]
    Výsledek není roven nule, takže přímky nemají stejný směr a jsou
    různoběžné.

    Současně platí
    \[
        ak-b
        =
        1\cdot2-2
        =
        0.
    \]
    Přímky jsou tedy kolmé.

    Průsečík určíme dosazením rovnice $y=2x$ do obecné rovnice:
    \begin{align*}
        x+2(2x)-5&=0,\\
        x+4x-5&=0,\\
        5x&=5,\\
        x&=1.
    \end{align*}

    Druhou souřadnici získáme ze směrnicové rovnice:
    \[
        y=2\cdot1=2.
    \]

    Průsečíkem je bod
    \[
        P[1;2].
    \]

    Platí tedy
    \[
        p\perp q
        \qquad\text{a}\qquad
        p\cap q=\{P[1;2]\}.
    \]
\end{example}

\begin{example}\label{exa:Mixed_line_forms_summary}
    Jsou dány tři přímky, z nichž je každá zadaná jiným způsobem:
    \[
        \begin{aligned}
            p:\quad x&=t,\\
                     y&=2t,
                     \qquad t\in\mathbb{R},
        \end{aligned}
    \]
    \[
        q:\quad 2x-y+5=0,
    \]
    \[
        r:\quad y=-\frac12x+\frac52.
    \]

    Určete vzájemnou polohu každé dvojice přímek. U různoběžných
    přímek určete také jejich průsečík.

    \begin{enumerate}[label=\alph*)]
        \item Přímky $p$ a $q$,
        \item přímky $p$ a $r$,
        \item přímky $q$ a $r$.
    \end{enumerate}

    Nejprve z jednotlivých zápisů získáme údaje popisující směr
    přímek.

    Přímka $p$ má směrový vektor
    \[
        \vec{u}_p=(1;2).
    \]

    Přímka $q$ má normálový vektor
    \[
        \vec{n}_q=(2;-1),
    \]
    a proto jako její směrový vektor můžeme zvolit
    \[
        \vec{u}_q=(1;2).
    \]

    Přímka $r$ má směrnici
    \[
        k_r=-\frac12.
    \]
    Jako její směrový vektor můžeme zvolit
    \[
        \vec{u}_r=(2;-1).
    \]

    \begin{enumerate}[label=\alph*), leftmargin=2em]

        \item \textbf{Přímky $p$ a $q$}

        Směrové vektory přímek jsou totožné:
        \[
            \vec{u}_p=\vec{u}_q=(1;2).
        \]
        Přímky tedy mají stejný směr.

        Bod $A[0;0]$ leží na přímce $p$. Dosadíme jej do obecné rovnice
        přímky $q$:
        \[
            2\cdot0-0+5=5\neq0.
        \]
        Bod $A$ na přímce $q$ neleží, a proto jsou přímky rovnoběžné
        a různé:
        \[
            p\parallel q.
        \]

        \item \textbf{Přímky $p$ a $r$}

        Vypočítáme skalární součin jejich směrových vektorů:
        \[
            \vec{u}_p\cdot\vec{u}_r
            =
            (1;2)\cdot(2;-1)
            =
            2-2
            =
            0.
        \]
        Přímky jsou tedy kolmé, a proto různoběžné.

        Průsečík určíme dosazením parametrického vyjádření přímky $p$
        do směrnicové rovnice přímky $r$:
        \begin{align*}
            2t
            &=
            -\frac12t+\frac52,\\
            \frac52t
            &=
            \frac52,\\
            t&=1.
        \end{align*}

        Pro $t=1$ dostáváme
        \[
            P[1;2].
        \]

        Platí tedy
        \[
            p\perp r
            \qquad\text{a}\qquad
            p\cap r=\{P[1;2]\}.
        \]

        \item \textbf{Přímky $q$ a $r$}

        Přímka $q$ má směrový vektor
        \[
            \vec{u}_q=(1;2),
        \]
        zatímco přímka $r$ má směrový vektor
        \[
            \vec{u}_r=(2;-1).
        \]
        Jejich skalární součin je
        \[
            (1;2)\cdot(2;-1)=2-2=0.
        \]
        Přímky jsou tedy také kolmé.

        Průsečík určíme dosazením směrnicové rovnice přímky $r$ do
        obecné rovnice přímky $q$:
        \begin{align*}
            2x-\left(-\frac12x+\frac52\right)+5&=0,\\
            2x+\frac12x-\frac52+5&=0,\\
            \frac52x+\frac52&=0,\\
            x&=-1.
        \end{align*}

        Druhou souřadnici dopočítáme ze směrnicové rovnice:
        \[
            y
            =
            -\frac12(-1)+\frac52
            =
            \frac12+\frac52
            =
            3.
        \]

        Průsečíkem je bod
        \[
            Q[-1;3].
        \]

        Platí tedy
        \[
            q\perp r
            \qquad\text{a}\qquad
            q\cap r=\{Q[-1;3]\}.
        \]

    \end{enumerate}

    Výsledky můžeme shrnout:
    \[
    \begin{array}{c|c|c}
        \text{dvojice přímek}
        &
        \text{vzájemná poloha}
        &
        \text{průsečík}
        \\
        \hline
        p,\ q
        &
        \text{rovnoběžné}
        &
        \varnothing
        \\[1mm]
        p,\ r
        &
        \text{kolmé}
        &
        P[1;2]
        \\[1mm]
        q,\ r
        &
        \text{kolmé}
        &
        Q[-1;3]
    \end{array}
    \]
\end{example}

Při určování vzájemné polohy přímek zadaných různými způsoby tedy není
nutné vždy provádět úplný převod. Často stačí z každého zápisu získat
směrový nebo normálový vektor, případně směrnici, a průsečík určit
přímým dosazením. Volba postupu závisí na konkrétním zadání a na tom,
který zápis vede k nejjednoduššímu výpočtu.
